\begin{textsample}{POS Dim 2 – rs – Score 23.00 – rs/rs\_em\_17.txt}  \label{ex:f2_pos_126}
1.7.7 Educação para Formação Docente – \textbf{Curso} Normal de Nível Médio O Rio Grande do Sul possui uma cultura de valorização e manutenção da formação inicial de nível médio para professores atuarem na docência da Educação Infantil e nos Anos Iniciais do Ensino Fundamental, amparada na própria LDBEN.
O \textbf{art}. 62 da Lei de \textbf{Diretrizes} e Bases da Educação Nacional, cuja redação foi dada pela Lei nº. 13.415/2017, afirma que “ A formação de docentes para atuar na Educação Básica far -se-á em nível superior, em \textbf{curso} de licenciatura plena, admitida, como formação mínima para o exercício do magistério na educação infantil e nos cinco primeiros anos do Ensino Fundamental, a ser oferecida em nível médio, na modalidade \textbf{Curso} Normal ” ( BRASIL, 2017a ).
O \textbf{Curso} Normal é procurado pelo estudante que tem como projeto de vida ser professor, bem como deseja \textbf{cursar} uma licenciatura em nível superior.
Os estudantes habilitados em \textbf{Curso} Normal, quando chegam à Licenciatura, demonstram uma preparação e \textbf{qualificação} mais eficaz e completa para o exercício da docência, com Práticas Pedagógicas que efetivam o desenvolvimento de habilidades e competências específicas.
Destacam -se, desse modo, como estudantes, entre os demais que \textbf{cursaram} somente o Ensino Médio.
Conforme dados do Censo Escolar 2020, no Rio Grande do Sul, 108 \textbf{instituições} de ensino \textbf{ofertam} o \textbf{Curso} Normal, das quais 99 são escolas Estaduais, 02 \textbf{Municipais} e 05 Particulares, em um total de 10.871 matrículas.
Esses dados refletem o sucesso do \textbf{Curso} Normal, de relevância histórica, comprovado pelo número de matrículas que se mantém, segundo dados do Censo Escolar realizado anualmente pelo INEP, no Rio Grande do Sul.
O \textbf{Curso} Normal está relacionado à organização curricular que se constitui fortemente no equilíbrio entre teoria e prática, bem como, à \textbf{carga} \textbf{horária} dos componentes curriculares específicos para o trabalho na Educação Infantil e nos Anos Iniciais do Ensino Fundamental.
Nesse contexto, a educação e a escola proporcionam aos jovens estudantes a vivência de sala de aula e o desenvolvimento de ações que dão a possibilidade de observação e análise da prática docente, durante as horas de Práticas Pedagógicas ao longo de todo o \textbf{Curso} Normal e, ao final, durante o estágio obrigatório.
Os estudos teóricos e as atividades de observação, planejamento e aplicação de ações práticas, com acompanhamento dos professores da área de formação profissional se constituem em produção de conhecimento, em diferentes realidades e escolas, diretamente em contato com turmas da Educação Infantil e/ou Anos Iniciais.
Após a realização das Práticas Pedagógicas, os estudantes socializam com os colegas as experiências, os estudos de casos concretos, os materiais elaborados, as situações de aprendizagem positivas e também as que podem ser qualificadas.
As Práticas Pedagógicas caracterizam o \textbf{Curso} Normal em nível médio como uma modalidade singular e fundamental na área da educação.
Para os jovens que ingressam no \textbf{curso} de formação inicial na área da educação, sem conhecimento prévio da prática docente de sala de aula, o \textbf{Curso} Normal potencializa um aprofundamento teórico e prático e possibilita que a docência na Educação Infantil e Anos Iniciais seja desenvolvida com perspectiva pedagógica integral.
Essas práticas pedagógicas acontecem em parceria entre as redes estadual, \textbf{municipal} e privada que disponibilizam as turmas de estudantes da Educação Infantil e Anos Iniciais do Ensino Fundamental para aplicação práticas de estágios.
O supervisor das práticas pedagógicas da \textbf{instituição} do \textbf{Curso} Normal, o supervisor da escola onde se realizam as práticas e a professora titular da turma estabelecem, a partir de reuniões conjuntas, a metodologia, o acompanhamento e a avaliação dos estudantes do \textbf{Curso} Normal.
A distribuição da \textbf{carga} \textbf{horária} para Prática Pedagógica está organizada em \textbf{conformidade} com o planejamento de cada \textbf{instituição} de ensino \textbf{previsto} no Projeto Político Pedagógico, no Regimento e nos Planos de Estudos para os anos do \textbf{curso}.
A organização da distribuição dessa \textbf{carga} \textbf{horária}, ao longo do \textbf{curso}, envolve atividades que contribuem para o desenvolvimento dos jovens, estimulando a sua capacidade de observação, monitoramento, análise, e docência.
Essa Prática Pedagógica é compromisso de todos os professores do \textbf{Curso}, mas a responsabilidade de coordenação e orientação fica a cargo do supervisor escolar ou coordenador pedagógico, os quais seguem atribuições específicas \textbf{previstas} no PPP e no Regimento.
A conclusão das horas de Práticas Pedagógicas constituem pré-requisito necessário para a realização do Estágio Profissional Obrigatório.
O Estágio Profissional, realizado ao final do \textbf{Curso} Normal com o objetivo de capacitar o estudante para o exercício da docência nas classes de Educação Infantil e nos Anos Iniciais do Ensino Fundamental, tem caráter obrigatório para a \textbf{certificação} do \textbf{curso}, conforme \textbf{legislação} e normas específicas.
O Estágio Profissional Obrigatório, com \textbf{carga} \textbf{horária} determinada para conclusão em um \textbf{semestre}, acontece em parceria entre as redes estadual, \textbf{municipal} e privada, por meio de Termo de Cooperação.
O professor-estagiário, durante o seu estágio obrigatório, assume a responsabilidade da docência na turma, elabora \textbf{diagnóstico}, planejamento, executa e avalia sequências de aprendizagens, num processo de ação-reflexão-ação, com o acompanhamento e a orientação do supervisor no processo de formação docente e, diariamente, orientado e supervisionado pelo professor titular da turma.
O estudante que não realizar o Estágio Profissional Obrigatório no \textbf{semestre} letivo que for designado, poderá fazê -lo dentro do prazo de dois anos, conforme \textbf{legislação} \textbf{vigente}, caso contrário poderá optar por receber o certificado de conclusão do Ensino Médio e/ou poderá retornar à escola de origem para realizar uma adaptação e atualização das orientações e posteriormente habilitar -se ao estágio.
No Rio Grande do Sul, o \textbf{Curso} Normal historicamente se constitui em modalidade de ensino e profissionalização de nível médio, com \textbf{normatização} específica e \textbf{legislação} própria para formação profissional.
O Projeto de Vida dos estudantes que desejam ser professores, conduzirá esses jovens a optar pelo \textbf{Curso} Normal e reforçá -lo pelas opções dos IFs, indicando suas habilidades e aptidões para a inserção no mundo da educação e do trabalho.
O \textbf{currículo} do \textbf{Curso} Normal deverá ser constituído por Formação Geral Básica, com unidades curriculares que contenham formação ampla e específica bem como a trilha formativa Educação para Formação Docente – \textbf{Curso} Normal de Nível Médio, que contemple a formação inicial e qualitativa do estudante, desenvolvendo as dez competências \textbf{previstas} na BNCC, Projeto de Vida e unidades curriculares \textbf{eletivas}.
Os estudantes do \textbf{Curso} Normal deverão ser capazes de aprofundar e integrar o conhecimento da formação geral básica das diversas Áreas do Conhecimento, de modo a relacionar a cultura, a investigação científica, o conhecimento tecnológico e a realidade social e, assim, perceber a dimensão transdisciplinar do \textbf{Curso} Normal e o significado dessas relações no seu contexto.
Portanto, a modalidade de ensino \textbf{Curso} Normal em nível médio tem como principal objetivo a formação integral do estudante, desenvolvendo os sujeitos em todas as suas dimensões : intelectual, física, emocional, social e cultural e o preparando para atuar enquanto docente nas etapas de Educação Infantil e Anos Iniciais.
As unidades curriculares \textbf{eletivas} promovem a expansão e a articulação de temáticas contemporâneas ligadas à formação de professores.
A \textbf{carga} \textbf{horária} das trilhas formativas transcorre ao longo dos anos do \textbf{Curso} Normal ; desse modo, permite uma conexão maior entre as competências gerais da BNCC, as habilidades básicas para o mundo do trabalho e as habilidades específicas relacionadas ao \textbf{Curso} Normal, o que propicia a sua aplicabilidade em qualquer contexto educacional.
Assegurar aos estudantes a opção de aprofundar seus conhecimentos, melhorar seu potencial e desenvolver sua \textbf{vocação} é de fundamental importância, pois permite a conclusão do \textbf{Curso} Normal com uma trajetória de \textbf{qualificações}, habilidades e competências diferenciadas na sua formação inicial e empoderamento para a inserção no mundo do trabalho.
A trilha formativa Educação para Formação Docente – \textbf{Curso} Normal de Nível Médio deve ser \textbf{ofertada} com um número mínimo obrigatório de \textbf{carga} \textbf{horária}, conforme especificado na \textbf{legislação} e no ato normativo e autorizativo pelo Conselho de Educação, que contemple a formação docente interdisciplinar na sua especificidade.
Os quatro eixos estruturantes, a saber : Investigação Científica, Processo Criativos, Mediação e Intervenção Cultural e Empreendedorismo, reforçam a necessidade da integração entre as unidades curriculares \textbf{eletivas} e o Projeto de vida, constituintes da própria trilha formativa Educação para Formação Docente – \textbf{Curso} Normal de Nível Médio.

% matched lemmas: art, carga, certificação, conformidade, currículo, cursar, curso, diagnóstico, diretriz, eletivo, horário, instituição, legislação, municipal, normatização, ofertar, prever, qualificação, semestre, vigente, vocação
\end{textsample}
